% test-shapepar.tex — shapepar (v2.2) test
% shapepar API:
%   \usepackage{shapepar}
%   \shapepar[<scale>]{<shape_spec>} Paragraph text \par
%     Scale auto-computed if omitted. Shape applied via \parshape.
%   \Shapepar[<scale>]{<shape_spec>} Paragraph text \par
%     Boxed version (no centering, for float contexts).
%   \cutout{r|l}(h_off, v_off)
%     Precedes \shapepar. Following paragraphs wrap around the shape
%     with \cutoutsep gap (default 12pt).
%   Predefined shapes + convenience macros:
%     \squareshape / \squarepar{text}
%     \circleshape / \circlepar{text}
%     \diamondshape / \diamondpar{text}
%     \heartshape / \heartpar{text}
%     \starshape / \starpar{text}
%     \nutshape / \nutpar{text}
%     \hexagonshape / \hexagonpar{text}
%   \rectangleshape{height}{width} — generates rectangle shape spec
%     NOTE: height/width must be plain numbers (shape units, not cm/pt)
%           or use with explicit scale parameter.
% Limitations (CRITICAL):
%   - parshape-based: CANNOT span page breaks
%   - NO \\ line breaks inside shaped paragraph
%   - NO \vspace, \centering, or any vertical material inside shaped par
%   - NO displayed math inside shaped paragraph
%   - NO list environments (itemize/enumerate) inside shaped paragraph
%   - \cutout only supports {l} or {r} — no centered cutout
%   - shaped paragraph is ONE paragraph (delimited by \par)
%   - \Shapepar requires text as argument (braces), \shapepar takes
%     text after the shape spec up to \par
\documentclass[11pt,a4paper]{article}
\usepackage[utf8]{inputenc}
\usepackage{graphicx}
\usepackage{shapepar}
\usepackage{lipsum}
\usepackage{multicol}

\title{shapepar (v2.2) Test}
\author{Programmer Agent}
\date{2026-05-16}

\setlength{\cutoutsep}{10pt}

\begin{document}
\maketitle

% ============================================================
% Test 1: Basic right-side cutout with \squarepar
% \cutout{r} places the shaped paragraph on the right margin.
% \squarepar is a convenience macro: \shapepar\squareshape{text}\par
% The text inside the shape flows in a square. Following paragraphs
% wrap around the square with \cutoutsep gap.
% NOTE: \vspace and \\ are FORBIDDEN inside shaped paragraphs.
% We cannot put images inside the shaped paragraph because
% \includegraphics requires vertical material handling that
% conflicts with shapepar's parshape mechanism.
% ============================================================
\section{Test 1: Basic right-side cutout}

\cutout{r}(0pt,0pt)\squarepar{%
  This is text flowing inside a square shape on the right side
  of the page. shapepar uses \texttt{\textbackslash parshape}
  to control line widths. The shape auto-scales to fit the text.
  Lorem ipsum dolor sit amet, consectetur adipiscing elit.
  Sed do eiusmod tempor incididunt ut labore et dolore magna
  aliqua. Ut enim ad minim veniam, quis nostrud exercitation.%
}\par

This is the first paragraph after the cutout. The text should wrap
around the square shape on the right side. The cutout mechanism
computes an exclusion zone around the shape and applies parshape
to subsequent paragraphs. \lipsum[1]

This is the second paragraph. Once the cutout zone ends, text
returns to full width. The cutout height equals the square's
vertical extent. \lipsum[2]

% ============================================================
% Test 2: Left-side cutout
% \cutout{l} places the shaped paragraph on the left margin.
% ============================================================
\section{Test 2: Left-side cutout}

\cutout{l}(0pt,0pt)\squarepar{%
  Text flowing inside a square on the left margin. The cutout
  mechanism adjusts subsequent paragraphs to avoid this zone.
  Lorem ipsum dolor sit amet, consectetur adipiscing elit sed
  do eiusmod tempor incididunt ut labore et dolore magna aliqua
  ut enim ad minim veniam quis nostrud exercitation ullamco.%
}\par

This paragraph should wrap around the square shape on the left
side. \lipsum[3]

Another paragraph that may still be in the cutout zone. \lipsum[4]

% ============================================================
% Test 3: Diamond shape with cutout
% Test different predefined shapes to verify they work with
% the cutout mechanism.
% ============================================================
\section{Test 3: Diamond shape cutout}

\cutout{r}(0pt,0pt)\diamondpar{%
  Text flowing inside a diamond shape on the right. The diamond
  tapers at top and bottom, so the cutout zone is narrower
  there. Lorem ipsum dolor sit amet, consectetur adipiscing
  elit sed do eiusmod tempor incididunt ut labore.%
}\par

This paragraph wraps around the diamond. The cutout zone follows
the diamond's contour, not a simple rectangle. \lipsum[5]

\newpage

% ============================================================
% Test 4: Tall shape near page break
% shapepar is parshape-based and CANNOT span page breaks.
% When a shaped paragraph is near a page boundary, inter-line
% penalties discourage but do not prevent breaks. If the shape
% does break, rendering is corrupted. We test with a long
% shaped paragraph to see behavior near page bottom.
% EXPECTED: N/A (design limitation — same as cutwin/picinpar).
% The shape will either fit entirely on one page or corrupt.
% ============================================================
\section{Test 4: Shape near page break}

\lipsum[6]

\vspace{4cm}

\cutout{r}(0pt,0pt)\squarepar{%
  This is a longer shaped paragraph near the page break. If there
  is insufficient room on the current page, the shape may be pushed
  to the next page entirely (best case) or may break across pages
  causing rendering corruption (worst case). Lorem ipsum dolor sit
  amet, consectetur adipiscing elit. Sed do eiusmod tempor incididunt
  ut labore et dolore magna aliqua. Ut enim ad minim veniam, quis
  nostrud exercitation ullamco laboris nisi ut aliquip ex ea commodo
  consequat. Duis aute irure dolor in reprehenderit.%
}\par

This paragraph follows the shaped cutout. If the shape was pushed
to the next page, this text renders at full width. \lipsum[7]

\newpage

% ============================================================
% Test 5: Itemize after cutout
% The cutout zone applies to ALL subsequent paragraphs via
% parshape. itemize uses its own list machinery. We test
% whether list items respect the parshape constraint.
% NOTE: itemize INSIDE the shaped paragraph is impossible.
% ============================================================
\section{Test 5: Itemize after cutout}

\cutout{r}(0pt,0pt)\squarepar{%
  Square shape with text. Lorem ipsum dolor sit amet,
  consectetur adipiscing elit sed do eiusmod tempor.%
}\par

\begin{itemize}
  \item First item after the cutout zone. If parshape is still
    active, this text should wrap around the shape on the right.
  \item Second item, also potentially wrapped.
  \item Third item, may have returned to full width.
  \item Fourth item, likely at full width.
  \item Fifth item, definitely at full width.
\end{itemize}

Trailing text after the itemize environment.

% ============================================================
% Test 6: multicol with cutout
% shapepar uses parshape which applies to the enclosing
% column's \hsize. Inside multicols the column width is
% narrower. This is a stress test — behavior may vary.
% ============================================================
\section{Test 6: multicol with cutout}

\begin{multicols}{2}

\cutout{r}(0pt,0pt)\squarepar{%
  Small square inside multicols. Lorem ipsum dolor sit amet.%
}\par

Text in the first column after the cutout. Since parshape uses
\texttt{\textbackslash hsize}, the cutout should adapt to the
narrower column. \lipsum[8]

More text. \lipsum[9]

\lipsum[10-11]

\end{multicols}

% ============================================================
% Test 7: Decorative diamond shape (package's intended use case)
% shapepar's primary design goal is decorative paragraph shapes,
% NOT figure wrapping. This test demonstrates the intended
% functionality: text flowing into interesting shapes.
% NOTE: \circlepar causes arithmetic overflow in some contexts;
% we use \diamondpar instead which is more numerically stable.
% ============================================================
\section{Test 7: Decorative diamond shape}

\diamondpar{%
  Lorem ipsum dolor sit amet, consectetur adipiscing elit.
  Sed do eiusmod tempor incididunt ut labore et dolore magna
  aliqua. Ut enim ad minim veniam, quis nostrud exercitation
  ullamco laboris nisi ut aliquip ex ea commodo consequat.
  Duis aute irure dolor in reprehenderit in voluptate velit
  esse cillum dolore eu fugiat nulla pariatur. Excepteur sint
  occaecat cupidatat non proident, sunt in culpa qui officia
  deserunt mollit anim id est laborum.%
}

\lipsum[12]

\end{document}
